\documentclass[12pt]{article}
\usepackage[a4paper,margin=1in]{geometry}
\usepackage[utf8]{inputenc}
\usepackage[T1]{fontenc}
\usepackage{amsmath,amssymb}
\usepackage{textcomp}
\usepackage{parskip}
\begin{document}
\begin{center}
{\Large \textbf{2018 STEP III}}
\end{center}

\subsection*{Page 1}
\begin{flushleft}
\bigskip\textbf{Section A: Pure Mathematics}\par
\bigskip\textbf{Question 1.} (i) The function f is given by\par
1 1\par
\[f(β) = β (β = 0) .\]
\[- β - β2 \leq\]
\[Find the stationary point of the curve y = f(β) and sketch the curve.\]
\[Sketch also the curve y = g(β) , where\]
3 1\par
\[g(β) = β + (β = 0) .\]
\[β - β2 \leq\]
\textbf{(ii)} Let u and v be the roots of the equation\par
\[x^2 + αx + β = 0 ,\]
1 1 1\par
\[where β = 0 . Obtain expressions in terms of α and β for u + v + and + + uv .\]
\[\leq uv u v\]
1 1 1\par
\textbf{(iii)} Given that u + v + = 1 , and that u and v are real, show that + + uv \leq 1 .\par
uv - u v -\par
\textbf{(iv)} Given instead that u + v + = 3 , and that u and v are real, find the greatest value\par
uv\par
1 1\par
of + + uv .\par
u v\par
\bigskip\textbf{Question 2.} The sequence of functions y , y , y , . . . is defined by y = 1 and, for n \geq 1 ,\par
0 1 2 0\par
1 dnz\par
\[y = ( 1)n ,\]
n - z dxn\par
\[where z = e\]
\[-x^2\]
.\par
dy\par
\textbf{(i)} Show that n = 2xy y for n \geq 1 .\par
n n+1\par
dx -\par
\textbf{(ii)} Prove by induction that, for n \geq 1 ,\par
\[y n+1 = 2xy n 2ny n-1 .\]
-\par
\[Deduce that, for n \geq 1 ,\]
\[y n 2 +1 - y n y n+2 = 2n y n 2 - y n-1 y n+1 ) + 2y n 2 .\]
\[\leq\]
\textbf{(iii)} Hence show that y n 2 - y n-1 y n+1 > 0 for n \geq 1 .\par
\end{flushleft}
\subsection*{Page 2}
\begin{flushleft}
\bigskip\textbf{Question 3.} Show that the second-order differential equation\par
\[x2y \leq\leq + (1 2p)x y \leq + (p^2 q^2) y = f(x) ,\]
- -\par
where p and q are constants, can be written in the form\par
\[xa xb(xcy) \leq \leq = f(x) , ( )\]
∗\par
\[\leq \geq\]
where a, b and c are constants.\par
\textbf{(i)} Use ( ) to derive the general solution of the equation\par
∗\par
\[x2y \leq\leq + (1 2p)xy \leq + (p^2 q^2)y = 0\]
- -\par
in the different cases that arise according to the values of p and q.\par
\textbf{(ii)} Use ( ) to derive the general solution of the equation\par
∗\par
\[x2y \leq\leq + (1 2p)xy \leq + p2y = xn\]
-\par
in the different cases that arise according to the values of p and n.\par
\bigskip\textbf{Question 4.} The point P (a sec θ, b tan θ) lies on the hyperbola\par
\[x^2 y^2\]
\[= 1 ,\]
\[a^2 - b^2\]
\[where a > b > 0 . Show that the equation of the tangent to the hyperbola at P can be\]
written as\par
\[bx ay sin θ = ab cos θ .\]
-\par
x y x y\par
\textbf{(i)} This tangent meets the lines = and = at S and T , respectively.\par
a b a - b\par
How is the mid-point of ST related to P ?\par
\textbf{(ii)} The point Q(a sec φ, b tan φ) also lies on the hyperbola and the tangents to the hyper-\par
bola at P and Q are perpendicular. These two tangents intersect at (x, y).\par
\[Obtain expressions for x^2 and y^2 in terms of a, θ and φ.\]
\[Hence, or otherwise, show that x^2 + y^2 = a^2 b^2.\]
-\par
\end{flushleft}
\subsection*{Page 3}
\begin{flushleft}
\bigskip\textbf{Question 5.} The real numbers a , a , a , . . . are all positive. For each positive integer n, A and G are\par
1 2 3 n n\par
defined by\par
a + a + + a\par
1 2 n 1/n\par
\[A = \cdots and G = a a a .\]
n n 1 2 n\par
n \cdots\par
\[\leq \geq\]
\textbf{(i)} Show that, for any given positive integer k,\par
\[(k + 1)(A G ) \geq k(A G )\]
k+1 k+1 k k\par
- -\par
if and only if\par
\[λk+1 (k + 1)λ + k \geq 0 ,\]
k - k\par
a k+1\par
\[where λ = k+1 .\]
k ( G )\par
k\par
\textbf{(ii)} Let\par
\[f(x) = xk+1 (k + 1)x + k ,\]
-\par
\[where x > 0 and k is a positive integer. Show that f(x) \geq 0 and that f(x) = 0 if and\]
\[only if x = 1 .\]
\textbf{(iii)} Deduce that:\par
\[(a) A \geq G for all n;\]
n n\par
\[(b) if A = G for some n, then a = a = = a .\]
n n 1 2 n\par
\cdots\par
\end{flushleft}
\subsection*{Page 4}
\begin{flushleft}
\bigskip\textbf{Question 6.} (i) The distinct points A, Q and C lie on a straight line in the Argand diagram, and\par
q a\par
represent the distinct complex numbers a, q and c, respectively. Show that - is\par
c a\par
\[real and hence that (c a)(q ∗ a ∗ ) = (c ∗ a ∗ )(q a) . -\]
- - - -\par
∗ ∗\par
\[Given that aa = cc = 1, show further that\]
∗\par
\[q + acq = a + c .\]
\textbf{(ii)} The distinct points A, B, C and D lie, in anticlockwise order, on the circle of unit\par
∗\par
\[radius with centre at the origin (so that, for example, aa = 1). The lines AC and BD\]
meet at Q. Show that\par
∗\par
\[(ac bd)q = (a + c) (b + d) ,\]
- -\par
where b and d are complex numbers represented by the points B and D respectively,\par
and show further that\par
∗\par
\[(ac bd)(q + q ) = (a b)(1 + cd) + (c d)(1 + ab) .\]
- - -\par
\textbf{(iii)} The lines AB and CD meet at P , which represents the complex number p. Given\par
\[that p is real, show that p(1 + ab) = a + b . Given further that ac bd = 0 , show that\]
\[- \leq\]
∗\par
\[p(q + q ) = 2 .\]
\end{flushleft}
\subsection*{Page 5}
\begin{flushleft}
\bigskip\textbf{Question 7.} (i) Use De Moivre's theorem to show that, if sin θ = 0 , then\par
\[\leq\]
(cot θ + i)2n+1 (cot θ i)2n+1 sin (2n + 1) θ\par
\[- - = ,\]
\[2i sin2n+1 θ\]
for any positive integer n.\par
Deduce that the solutions of the equation\par
2n + 1 2n + 1\par
\[xn xn-1 + + ( 1)n = 0\]
( 1 ) - ( 3 ) \cdots -\par
are\par
\[m\pi\]
\[x = cot2\]
( 2n + 1 )\par
\[where m = 1, 2, . . . , n .\]
\textbf{(ii)} Hence show that\par
n\par
\[m\pi n (2n 1)\]
\[cot2 = - .\]
( 2n + 1 ) 3\par
\[m\int=1\]
\textbf{(iii)} Given that 0 < sin θ < θ < tan θ for 0 < θ < 1 \pi, show that\par
\[cot2 θ < < 1 + cot2 θ.\]
θ2\par
Hence show that\par
\infty\par
\[1 \pi2\]
\[= .\]
\[m^2 6\]
\[m\int=1\]
\bigskip\textbf{Question 8.} In this question, you should ignore issues of convergence.\par
\textbf{(i)} Let\par
1 f(x -1)\par
\[I = dx ,\]
\[\sum 1 + x\]
where f(x) is a function for which the integral exists.\par
Show that\par
\infty\par
n+1 f(y)\par
\[I = dy\]
\[\sum y(1 + y)\]
\[n\int=1 n\]
\[and deduce that, if f(x) = f(x + 1) for all x, then\]
1 f(x)\par
\[I = dx .\]
\[\sum 1 + x\]
\textbf{(ii)} The fractional part, x , of a real number x is defined to be x x where x is the\par
\{ \} - ( ) ( )\par
\[largest integer less than or equal to x. For example 3.2 = 0.2 and 3 = 0 .\]
\{ \} \{ \}\par
Use the result of part (i) to evaluate\par
1 x -1 1 2x -1\par
\{ \} dx and \{ \} dx .\par
\[\sum 1 + x \sum 1 + x\]
0 0\par
\end{flushleft}
\subsection*{Page 6}
\begin{flushleft}
\bigskip\textbf{Section B: Mechanics}\par
\bigskip\textbf{Question 9.} A particle P of mass m is projected with speed u along a smooth horizontal floor directly\par
towards a wall. It collides with a particle Q of mass km which is moving directly away from\par
the wall with speed v . In the subsequent motion, Q collides alternately with the wall and\par
with P . The coefficient of restitution between Q and P is e, and the coefficient of restitution\par
between Q and the wall is 1.\par
Let u and v be the velocities of P and Q, respectively, towards the wall after the nth\par
n n\par
collision between P and Q.\par
\textbf{(i)} Show that, for n \geq 2,\par
\[(1 + k)u n (1 k)(1 + e)u n-1 + e(1 + k)u n-2 = 0 . ( )\]
- - ∗\par
\textbf{(ii)} You are now given that e = 1 and k = 1 , and that the solution of ( ) is of the form\par
2 34 ∗\par
\[u = A 7 n + B 5 n (n \geq 0) ,\]
n 10 7\par
\[\leq \geq \leq \geq\]
where A and B are independent of n. Find expressions for A and B in terms of u\par
and v .\par
\[Show that, if 0 < 6u < v , then u will be negative for large n.\]
0 0 n\par
\bigskip\textbf{Question 10.} A uniform disc with centre O and radius a is suspended from a point A on its circumference,\par
so that it can swing freely about a horizontal axis L through A. The plane of the disc is\par
perpendicular to L. A particle P is attached to a point on the circumference of the disc.\par
The mass of the disc is M and the mass of the particle is m.\par
In equilibrium, the disc hangs with OP horizontal, and the angle between AO and the\par
downward vertical through A is β. Find sin β in terms of M and m and show that\par
AP 2M\par
\[= .\]
a ( M + m\par
The disc is rotated about L and then released. At later time t, the angle between OP and\par
the horizontal is θ; when P is higher than O, θ is positive and when P is lower than O, θ is\par
negative. Show that\par
\[1 Iθ˙2 + (1 sin β)ma2θ˙2 + (m + M )ga cos β (1 cos θ)\]
2 - -\par
is constant during the motion, where I is the moment of inertia of the disc about L.\par
\[Given that m = 3 M and that I = 3 M a^2, show that the period of small oscillations is\]
2 2\par
3a\par
\[3\pi .\]
( 5g\par
\end{flushleft}
\subsection*{Page 7}
\begin{flushleft}
\bigskip\textbf{Question 11.} A particle is attached to one end of a light inextensible string of length b. The other end of\par
the string is attached to a fixed point O. Initially the particle hangs vertically below O. The\par
particle then receives a horizontal impulse.\par
The particle moves in a circular arc with the string taut until the acute angle between the\par
string and the upward vertical is α, at which time it becomes slack. Express V , the speed of\par
the particle when the string becomes slack, in terms of b, g and α.\par
Show that the string becomes taut again a time T later, where\par
\[gT = 4V sin α ,\]
and that just before this time the trajectory of the particle makes an angle β with the\par
\[horizontal where tan β = 3 tan α .\]
When the string becomes taut, the momentum of the particle in the direction of the string\par
is destroyed. Show that the particle comes instantaneously to rest at this time if and only if\par
\[1 + \sqrt\{3\}\]
\[sin2 α = .\]
\end{flushleft}
\subsection*{Page 8}
\begin{flushleft}
\bigskip\textbf{Section C: Probability and Statistics}\par
\bigskip\textbf{Question 12.} A random process generates, independently, n numbers each of which is drawn from a uniform\par
(rectangular) distribution on the interval 0 to 1. The random variable Y is defined to be\par
k\par
the kth smallest number (so there are k 1 smaller numbers).\par
-\par
\textbf{(i)} Show that, for 0 \leq y \leq 1 ,\par
n\par
n\par
\[P Y \leq y) = ym (1 y)n-m . ( )\]
k\par
(m) - ∗\par
\[\leq m\int=k\]
\textbf{(ii)} Show that\par
n n 1\par
\[m = n -\]
(m) (m 1)\par
-\par
n\par
and obtain a similar expression for (n m) .\par
- (m)\par
Starting from ( ), show that the probability density function of Y is\par
k\par
∗\par
n\par
n\par
-\par
1 yk-1\par
(1\par
y)n-k\par
.\par
(k 1) -\par
-\par
Deduce an expression for\par
yk-1(1 y)n-k\par
dy .\par
\[\sum -\]
\textbf{(iii)} Find E(Y ) in terms of n and k.\par
k\par
\end{flushleft}
\subsection*{Page 9}
\begin{flushleft}
\bigskip\textbf{Question 13.} The random variable X takes only non-negative integer values and has probability generating\par
function G(t). Show that\par
\[P(X = 0 or 2 or 4 or 6 . . .) = G (1) + G ( 1) .\]
2 -\par
\[\leq \geq\]
You are now given that X has a Poisson distribution with mean λ. Show that\par
\[G(t) = e\]
-λ(1-t)\par
.\par
\textbf{(i)} The random variable Y is defined by\par
\[kP(X = r) if r = 0, 2, 4, 6, . . . ,\]
\[P(Y = r) = ⎧\]
⎨0 otherwise,\par
⎩\par
where k is an appropriate constant.\par
cosh λt\par
Show that the probability generating function of Y is .\par
cosh λ\par
\[Deduce that E(Y ) < λ for λ > 0 .\]
\textbf{(ii)} The random variable Z is defined by\par
\[cP(X = r) if r = 0, 4, 8, 12, . . . ,\]
\[P(Z = r) = ⎧\]
⎨0 otherwise,\par
⎩\par
where c is an appropriate constant.\par
\[Is E(Z) < λ for all positive values of λ ?\]
\end{flushleft}
\end{document}
